% document formatting
\documentclass[10pt]{article}
\usepackage[utf8]{inputenc}
\usepackage[left=1in,right=1in,top=1in,bottom=1in]{geometry}
\usepackage[T1]{fontenc}
\usepackage{xcolor}

% math symbols, etc.
\usepackage{amsmath, amsfonts, amssymb, amsthm}

% lists
\usepackage{enumerate}

% images
\usepackage{graphicx} % for images
\usepackage{tikz}

% code blocks
% \usepackage{minted, listings} 

% verbatim greek
\usepackage{alphabeta}

\graphicspath{{./assets/images/Week 1}}

\title{02-614 Week 1 \\ \large{String Algorithms}}
 
\author{Aidan Jan}

\date{\today}

\begin{document}
\maketitle

\section*{Strings}
A string is a data structure that represents a sequence of letters drawn from an alphabet.  Some examples:
\begin{itemize}
	\item ``abcdefg''
	\item ``hello''
	\item ``ACCGTACGG\dots''
\end{itemize}
This course covers algorithms used for processing strings.

\section*{Exact Matching}
One of the first questions we might ask is where a string exists in another string, if it exists.  Formally, given a string $T$ and a string $P$, report all places in $T$ where $P$ occurs.

\subsection*{Na\"ive Algorithm}
\begin{verbatim}
for every pos i in T:
    for every pos j in P:
        check if T[i + j] = P[j]
\end{verbatim}
This algorithm essentially checks every position in the string $T$ to see if $P$ begins at that position.  This has a running time of $O(|T| \times |P|)$, which is slow.  Ideally, we want an algorithm with $O(|T| + |P|) = O(|T|)$ instead, since that is theoretically the fastest time possible.

\subsection*{The $Z$ Algorithm}
This algorithm does exact matching in linear time.\\\\
Define $Z$-values.  $Z_i(P) :=$ the length of the longest substring starting at $i$ in $P$ that matches a prefix of $P$.  (e.g., a repeated substring in $P$).
\begin{verbatim}
def Zmatch(T, P):
    S = concat(P, '$', T)
    print all i for which Z_i = |P|
\end{verbatim}
Essentially, if we knew all the $Z_i$'s, then finding an exact matching is trivial.  The only issue now is, how do we compute the $Z_i$'s?

\subsubsection*{Computing $Z_i$ Values}
Assume we have computed $Z_2, \dots, Z_{k - 1}$, and $\ell$ and $r$, where $\ell, r :=$ the left and right edges of the rightmost non-zero $z_i$ values computed so far.  The section of $z$-values between $\ell$ and $r$ are called the $z$-box.  From here we have a few cases:
\begin{enumerate}
    \item Case $k > r$:  The $Z$-box ends before index $k$.  If this case occurs, we have to brute-force compute $Z_k$
    \item Case $k \leq r$:  The $Z$-box ends after index $k$.  We know that $\ell < k$ since the start of a $z$-box must be within the part we already computed.  Therefore $\ell$ is at most $k - 1$.  We know for a fact that the characters in a $Z$-box matches the characters at the beginning of the string.  Therefore, since $k$ is in our $Z$-box, there must exist a $k'$ in the same box at the beginning at the string which we have already calculated.  Now, define $\beta$ as the section of the string between $k$ and $r$.  Now we have two subcases:
    \begin{enumerate}
	    \item Subcase $Z_{k'} < \beta$:  We know that $Z_k = Z_{k'}$.
	    \item Subcase $Z_{k'} \geq \beta$:  If this is true, then $Z_{k'}$ extends past the end of the $Z$-box.  Therefore, we skip to $r$, and compute from that point using brute force.
    \end{enumerate}
\end{enumerate}


\end{document}